\documentclass[10pt]{article}
\usepackage{amsmath, amssymb, amsthm}
\usepackage[mathscr]{euscript}
\usepackage{setspace}
\usepackage[pdftex]{graphicx}
\usepackage{booktabs}
\usepackage[dvipsnames]{xcolor}
\usepackage{minted}


% Setting Margins
\setlength{\textheight}{9in} \setlength{\topmargin}{-.5in}
\setlength{\textwidth}{6.5in} \setlength{\oddsidemargin}{0in}
\setlength{\evensidemargin}{0in}

\title{COMP152 --- SP25 --- Lab 5 \\ Loop and Big-O Worksheet}
\author{}
\begin{document}
\maketitle

\thispagestyle{empty}
\subsection*{Loops: While and For}

For each loop shown below:
\begin{itemize}
    \item Determine \textit{exactly} (either a literal value or a function of $n$) how many times the $\ldots$ line will be executed.
    \item Determine the Big-0 characterization of the loop. 
    \item Translate to the equivalent while/for loop. This means a loop with a counter (i,j,k, etc.) that takes on the exact same sequence of values, in the same order, repeats the same number of times, and has the same Big-O characterization.
\end{itemize}

\begin{enumerate}
    \item 
    \begin{minted}{python3}
    j = 2
    while j < n:
        ...
        j += 3
    \end{minted}
    \vspace{2in}
    \item
    \begin{minted}{python3}
    for i in range(0,n-1,n//3):
        ...
    \end{minted}
    \newpage \thispagestyle{empty}
    \item
    \begin{minted}{python3}
    k = n+5
    while k >= 4:
        ...
        k -= 2
    \end{minted}
    \vspace{2in}
    \item
    \begin{minted}{python3}
    for i in range(n):
        for j in range(n-1,i-1,-1):
            ...
    \end{minted}
    \vspace{3in}
    \item 
    \begin{minted}{python}
       i = 0
       while i<=(n-1):
          for j in range(1,n,2):
             ...
        i += 2
    \end{minted}
\end{enumerate}
\newpage \thispagestyle{empty}
\subsection*{Big-O Analysis}

Determine the Big-0 characterization of the worst case time for each of the following. Clearly identify/name the $n$ for each. 

\begin{enumerate}
    \item \begin{minted}{python3}
        def foo(l : List[int]):
            s = 0
            for i in range(2**10):
                k = 1
                while k < len(l):
                    s += l[k]
                    k *= 2
            return s
    \end{minted}
    \vspace{1.5in}

    \item \begin{minted}{python3}
        def foo(l : List[int], a : int, b : int):
            k = []
            for i in range(len(l)):
                if a < l[i] < b:
                    k.append(l[i])
            return k
    \end{minted}
    \vspace{1.5in}

    \item \begin{minted}{python3}
        def foo(l : List[int],a : int, b : int):
            i = 0
            while i < len(l):
               if a < l[i] < b:
                  del l[i]
                else:
                  i += 1
            return 
    \end{minted}
    \newpage \thispagestyle{empty}

    \item \begin{minted}{python3}
        def foo(l : List[int], k : int):
            l.sort()
            return l[len(s)-k:] 
    \end{minted}
    \vspace{2in}
    \item \begin{minted}{python3}
        def foo(l : List[int]):
            for i in range(1,n):
                j = i
                while j>0 and l[j] < l[j-1]:
                    l[j],l[j-1] = l[j-1],l[j]
                    j -= 1
            return 
    \end{minted}
    \vspace{2in}
    \item \begin{minted}{python3}
        def foo(l : List[int]):
            i = 1
            while i < len(l):
                for j in range(0,len(l),i):
                    l[j] = l[j]+l[j+2*i-1]
                i *= 2
            return         
    \end{minted}
\end{enumerate}

\end{document}
